\documentclass[12pt]{article}
\usepackage[margin=1in, bottom = 1.5in]{geometry} 
\usepackage{amsmath,amsthm,amssymb,amsfonts, enumitem, float, fancyhdr, color, comment, graphicx, environ,setspace, subcaption, lastpage, tabularx, booktabs, caption}
\usepackage{tikz, tikz-cd}
\captionsetup[subfigure]{width=0.85\linewidth}
\pagestyle{fancy}
\fancyhf{}
\lhead{Galois Theory for Rats}  
\rhead{Galois Theory 101}
\fancyfoot[R]{Page \thepage\ of \pageref{LastPage}}
\newlength{\myfootheight}
\setlength{\myfootheight}{35pt} 
\renewcommand{\footrulewidth}{0.4pt}
\setlength{\headheight}{35pt}
\pagenumbering{arabic}


% Define the theorem, lemma, proposition and corollary environmentsc

\newtheorem{filler}{Filler}
\theoremstyle{definition}
\newtheorem{corollary}{Corollary}[section]
\theoremstyle{definition}
\newtheorem{definition}{Definition}[section]
\theoremstyle{definition}
\newtheorem{example}{Example}[section]
\theoremstyle{definition}
\newtheorem{theorem}{Theorem}[section]
\theoremstyle{definition}
\newtheorem{lemma}{Lemma}[section]
\theoremstyle{definition}
\newtheorem{proposition}{Proposition}[section]
\theoremstyle{definition}

\AtBeginDocument{%
  \mathchardef\stdcomma=\mathcode`,
  \mathcode`,="8000
}
\begingroup\lccode`~=`, \lowercase{\endgroup\def~}{\stdcomma\,}

\title{Galois Theory for Rats - Galois Theory 101}
\author{Skyler Hu}

\begin{document}
\linespread{1.05}
\maketitle
\newpage
\tableofcontents
\newpage

\begin{section}{Introduction}
    After our prerequisites (Field Extensions, Field Representation \& Construction), hopefully you have gained some basic ideas of 
    field theory. If so, congratulations! You are officially a graduate of Galois Theory kindergarten and can now move on to the  
    \textit{real} Galois Theory. Warning: the math is going to get insidious from now on, so be prepared.\\
    \break
    In this fourth handout of Galois Theory for Rats, we take the crucial step from fields to groups -- the core of Galois Theory, and 
    show how groups, fields and polynomials are closely interrelated. After some preliminary definitions, we get to enjoy the \textit{Galois 
    Correspondence} (or \textit{The Fundamental Theorem of Galois Theory}), where everything we've learned will start to fall together and 
    (ideally) you will understand why Galois Theory is considered one of the most beautiful branches of mathematics. Note that from now on, 
    logic and proofs take up a significantly greater part of the handouts, and unlike the previous ones where proofs for most results and optional,
    you will be forced to understand proofs, some of them rather painful. You are encouraged to read what you don't understand as many times 
    as you need to -- this is the journey of every junior mathematician.
\end{section}

\newpage

\begin{section}{Fields' Groups and Groups' Fields}
    In this section we introduce some definitions central to Galois Theory. These definitions are mainly established through the process of 
    creating groups from fields, fields from groups and both from polynomials. We start with another definition ending in `-morphism':
    \begin{definition}{(Field Automorphism)}
        An isomorphism \(\sigma\) of field \(K\) is called an \textit{autormorphism} if it maps every element in \(K\) onto elements in \(K\).
        The collection of automorphisms of a field \(K\) is denoted \(Aut(K)\).
    \end{definition}
    There exists many automorphisms of a field \(K\) since there are multiple ways to shuffle the elements of \(K\) around, and it follows naturally 
    that we'd be interested in a special category of field automorphisms: these automorphisms map some elements onto themselves.
    \begin{definition}
        An automorphism \(\sigma\in Aut(K)\) is said to \textit{fix} an element \(\alpha \in K\) if \(\sigma \alpha = \alpha\). If \(F\) is a subset/subfield 
        of \(K\), then an automorphism of \(K\) is said to \textit{fix \(F\)} if it fixes every element in \(F\). \(Aut(K/F)\) is defined as the collection 
        of automorphisms of \(K\) fixing \(F\).
    \end{definition}

    \begin{proposition}{(Automorphisms form a group.)}
        The collection of automorphisms of field \(K\), \(Aut(K)\), is a group under composition. It has subgroup \(Aut(K/F)\) for \(F\) a subfield of \(K\).
        \begin{proof}
            Can you see that this is obvious? (Hint: check group axioms)
        \end{proof}
    \end{proposition}

    Now that we've seen some very basic properties of automorphisms, we're ready to address a common statement made about Galois Theory: the permutation of roots. 
    Though at this point we still have a long way to go until we learn precisely how root permutation relates to solvability, right now we can show that automorphisms 
    permute the roots of polynomials. More specifically:

    \begin{proposition}{(Automorphism permutes roots.)}
        Let \(K/F\) be a field extension and let \(\alpha\) be algebraic over \(F\). Then for any automorphism \(\sigma \in Aut(K/F)\), \(\sigma\alpha\) is a root 
        of the minimal polynomial of \(\alpha\) over \(F\), i.e. \(Aut(K/F)\) permutes the roots of irreducible polynomials, and any polynomial in \(F[x]\) having 
        \(\alpha\) as a root also has \(\sigma \alpha\) as a root.
        \begin{proof}{(Short and Easy!)}
            Let \(\alpha\) be a root of the polynomial \(f(x)\in F[x]\)
            \[f(x) = a_n x^n a_{n-1} x^{n-1} + \ldots + a_1 x + a_0.\]
            Then \(a_n \alpha^n + a_{n-1} \alpha^{n-1} + \ldots + a_1 \alpha + a_0=0.\) Applying \(\sigma\) to each term gives us 
            \[\sigma(a_n \alpha^n) + \sigma(a_{n-1} \alpha^{n-1}) + \ldots + \sigma(a_1 \alpha) + \sigma(a_0) = \sigma(0).\]
        \end{proof}
        \noindent\textit{Note.} The \textit{minimal polynomial} of \(\alpha\) over \(F\) is defined to be the irreducible polynomial \(m(x)\in F[x]\) of least degree
        having \(\alpha\) as a root; in a sense, the most simple polynomial with \(\alpha\) as a root. It is often algebraic shorthand to say that something (in this case \(\sigma\alpha\))
        is the root of \(m(x)\) instead of `every polynomial having \(\alpha\) as a root also has \(\sigma\alpha\) as a root.'
    \end{proposition}

    \begin{example}
        \begin{enumerate}
            \item Consider the field extension \(\mathbb{Q}(\sqrt{2})\). If \(\tau\) is an automorphism of \(\mathbb{Q}(\sqrt{2})\) (i.e. \(\tau\in Aut(\mathbb{Q}(\sqrt{2}))\)),
            then \(\tau(\sqrt{2}) = \pm \sqrt{2}\), which are both roots of the minimal polynomial of \(\sqrt{2}\) over \(\mathbb{Q}\), which is \(x^2-2\).\\
            Since \(\tau\) fixes \(\mathbb{Q}\), we know that \(\tau\) is of the form
            \[\tau(a+b\sqrt{2}) = a\pm b\sqrt{2}.\]
            \(\tau(\sqrt{2}) = \sqrt{2}\) is the identity automorphism, and  \(\tau(\sqrt{2}) = -\sqrt{2}\) has order 2, so the group \(Aut(\mathbb{Q}(\sqrt{2})/\mathbb{Q})\)
            is the cyclic group of order 2.

            \item Consider the field extensions \(\mathbb{Q}(\sqrt[3]{2})/\mathbb{Q}\). You might be tempted to think that the automorphism group with \(\mathbb{Q}\) fixed is a cyclic group of order 3, but consider this:
            The minimal polynomial of \(\sqrt[3]{2}\) is \(x^3-2\), which has one root \(\sqrt[3]{2}\) with multiplicity 3. Therefore the automorphism fixing \(\mathbb{Q}\) also fixes \(\sqrt[3]{2}\) (and fixes 
            \(\sqrt[3]{4}\) for that matter), and \(Aut(\mathbb{Q}(\sqrt[3]{2})) = \{1\}/\)
        \end{enumerate}
    \end{example}

    We have seen that field extensions can create groups by considering the automorphism group of a field \(F\) (usually fixing a subfield of \(F\)). Now, in order to fully appreciate the relationship between these 
    algebraic structures, we will see how groups can create fields:

    \begin{proposition}{(Field Extensions from Groups)}
        Let \(H \leq Aut(K)\) be a subgroup of \(Aut(K)\). The collection \(F\) of elements in \(K\) that are fixed by all elements in \(H\) form a subfield of \(K\).
        \begin{proof}
            Check field axioms. Be sure to use the definition of an isomorphism that \(\tau(ab) = \tau(a)\tau(b)\).
        \end{proof}
    \end{proposition}

    \textbf{Proposition 2.3} gives us a field from considering an automorphism group: the elements that are fixed by every possible automorphism form a field. This leads naturally to the definition of a 
    \textit{fixed field}:

    \begin{definition}{(Fixed Field)}
        If \(H\) is a subgroup of \(Aut(K)\), the group of automorphisms of \(K\), then the subfield \(F\) of \(K\) that is fixed by all elements of \(H\) is called the fixed field of \(H\).
    \end{definition}

    \begin{example}{(Fixed fields)}
        We will use the same examples as \textbf{Example 2.1} about automorphism groups:
        \begin{enumerate}
            \item The fixed field of \(Aut(\mathbb{Q}(\sqrt{2})/\mathbb{Q})\) is \(\mathbb{Q}\). A fixed field must have every element fixed by 
            every possible automorphism.
            \item The fixed field of \(Aut(\mathbb{Q}(\sqrt[3]{2})/\mathbb{Q})\) is \(Aut(\mathbb{Q}(\sqrt[3]{2}))\) since the only automorphism is the identity.
        \end{enumerate}
    \end{example}

    In this section so far we have defined all the tools we need to stitch groups and fields together -- field create groups by automorphism, and groups create fields by 
    considering fixed fields. Keep these definitions in mind: we are accelerating!
\end{section}

\newpage

\begin{section}{Getting Galois}
    In this section we prove some results concerning the relationship between groups and fields using the definitions from Section 2. `A few results' may be an understatement -- the upcoming propositions 
    are central to piecing together the proof of the Galois Correspondence, where everything will start making sense. So fear not if the motivation behind these results seem unclear on a first read: 
    come back when you completely understand everything in Section 4 and I promise you it's going to look a lot prettier. (At least, I hope that you don't read this handout only once.)\\

    \textit{Note on Notation.} In order to keep the proofs short I have chosen not to define every single algebraic object in this section (and in the following sections). A general rule of thumb 
    is that \(E,G,H\) are groups and \(F,K,L\) are fields unless otherwise specified. We're also going to use \(E \leq F\) to denote `\(E\) is a subgroup of \(F\)' and \(F\subseteq K\) to denote 
    `\(F\) is a subfield of \(K\)',so apologies if you grew up using another notation.

    \begin{proposition}{(Inclusion-reversing Correspondence)}
        \begin{enumerate}
            \item If \(F_1 \subseteq F_2 \subseteq K\), then \(Aut(K/F_2) \leq Aut(K/F_1)\).
            \item If \(H_1 \leq H_2 \leq Aut(K)\) and \(H_1, H_2\) have fixed fields \(F_1, F_2\), then \(F_2 \subseteq F_1\).
        \end{enumerate}

        \begin{proof}
            \begin{enumerate}
                \item For every \(h\in Aut(K/F_2)\) fixing \(F_2\), \(h\) is also in \(Aut(K/F_1)\) since \(F_1 \subseteq F_2\), so \(Aut(K/F_2 \leq Aut(K/F_1))\).
                \item For every element fixed by some \(h\in H_2\), the element must be fixed by every \(h\in H_1\) since \(H_1 \leq H_2\). Therefore \(F_2 \subseteq F_1\). 
            \end{enumerate}
        \end{proof}
    \end{proposition}

    \begin{proposition}{(Degree of Extension vs. Order of Group)}
        Let \(E\) be the splitting field over \(F\) of the polynomial \(f(x) \in F[x]\). Then 
        \[\left|Aut(E/F)\right| \leq [E:F].\]
    \end{proposition}
    This proof requires a slightly more nuanced argument and is therefore delegated to the Proofs section. (If you're interested: this is one of the few instances 
    in algebra that proof by induction makes an appearance.)\\
    \break
    \textbf{Proposition 3.2} allows us to relate automorphism groups and fixed fields on a numerical level. Now we have the ideas we need to start looking at things
    named after Galois:

    \begin{definition}{(Galois group.)}
        If \(f(x)\) is a separable polynomial over \(F\), then the \textit{Galois group} of \(f(x)\) over \(F\)
        is the group of automorphisms of the splitting field of \(f(x)\) over \(F\). The Galois group of a field extension is denoted 
        \(Gal(K/F)\).
    \end{definition}

    If you're wondering where the separability condition comes from:

    \begin{definition}{(Galois extension.)}
        A field extension \(K/F\) is defined to be Galois if it is:
        \begin{enumerate}
            \item \textbf{(1)} Normal: if \(x\in K\) has minimial polynomial \(m(x) \in F[x]\), then every root of \(m(x)\) is in \(K\).
            \item \textbf{(2)} Separable: if \(x\in K\) has minimial polynomial \(m(x) \in F[x]\), then \(m(x)\) has distinct roots in its splitting field.
        \end{enumerate}
        Relationship between Galois group and Galois extension: the automorphism group of \(K/F\) is Galois if and only if \(K/F\) is a Galois extension. This 
        is by definition.
    \end{definition}

    \textit{Note.} We've introduced the definition of the normality and separability of field extensions here. If this bothers you, simply acknowledge that in 
    Galois Theory we usually only care about separable, irreducible (and sometimes minimal) polynomials. If you're wondering why a Galois extension is defined this way,
    it's a little secret :)\\

    
    Things start to get slightly abstract around here so we introduce some examples. Not everything may be spelled out from now on, so if you don't understand 
    where a statement came from you can think about it in terms of definitions.
    \begin{example}{(Galois Groups \& Extensions)}
        \begin{enumerate}
            \item The extension \(\mathbb{Q}(\sqrt{2})/\mathbb{Q}\) is Galois (why?) with Galois group isomorphic to \(\mathbb{Z}/2\mathbb{Z}\):
            \begin{align*}
                1 &: a+b\sqrt{2} \mapsto a+b\sqrt{2}\\
                \sigma &: a+b\sqrt{2} \mapsto a-b\sqrt{2}
            \end{align*}

            \item \(\mathbb{Q(\sqrt[3]{2})}/\mathbb{Q}\) is not Galois since \(\left|Aut(\mathbb{Q}(\sqrt[3]{2})/\mathbb{Q})\right| = 1\). \(\mathbb{Q}(\sqrt[3]{2})\)
            is also not the splitting field of any polynomial. (If you're considering \(x^3-2\): it has splitting field \(\mathbb{Q}(\sqrt[3]{2}, \omega)\) where \(\omega\)
            is a primitive third root of unity. More on roots of unity in our next handout.)

            \item \(\mathbb{Q}(\sqrt{2},\sqrt{3})/\mathbb{Q}\) is Galois since it is the splitting field of the separable polynomial \((x^2-2)(x^2-3)\). The elements in its 
            automorphism group (a Galois group) are:
            \begin{align*}
                1 &: \sqrt{2} \mapsto \sqrt{2}, \sqrt{3} \mapsto \sqrt{3}\\
                \sigma &: \sqrt{2} \mapsto -\sqrt{2}, \sqrt{3} \mapsto \sqrt{3}\\
                \tau &: \sqrt{2} \mapsto \sqrt{2}, \sqrt{3} \mapsto -\sqrt{3}\\
                \sigma\tau &: \sqrt{2} \mapsto -\sqrt{2}, \sqrt{3} \mapsto -\sqrt{3}
            \end{align*}
            which gives the automorphism group as the Klein-4 group.
        \end{enumerate}
    \end{example}
\end{section}

\newpage

\begin{section}{The Galois Correspondence}
    I will give you a peek:
    \begin{figure}[H]
        \centering
        \begin{tikzcd}[row sep = small, column sep = large, arrows=no head]
            K\arrow[d] \arrow[r]& 1\arrow[d]\\
            F_1\arrow[d] \arrow[r] & G_i\arrow[d]\\
            F_2\arrow[d] \arrow[r] & G_{i-1}\arrow[d]\\
            \ldots\arrow[d]\arrow[r] & \ldots\arrow[d]\\
            F_i\arrow[d]\arrow[r] & G_1\arrow[d]\\
            F\arrow[r] & H
        \end{tikzcd}
    \end{figure}

    No reaction? Here's the complete version:
    \begin{theorem}{(Fundamental Theorem of Galois Theory)}
        Let \(K/F\) be a Galois extension and set \(G=Gal(K/F)\). Then there is a bijection:
        \begin{figure}[H]
        \centering
        \begin{tikzcd}[row sep = small, column sep = large, arrows=no head]
            K\arrow[d] \arrow[r]& 1\arrow[d]\\
            F_1\arrow[d] \arrow[r] & H_i\arrow[d]\\
            \ldots\arrow[d]\arrow[r] & \ldots\arrow[d]\\
            F\arrow[r] & G
        \end{tikzcd}
        \caption{subfields of \(K \mapsto\) subgroups of \(G\) in reverse-inclusion}
    \end{figure}
    given by 
    \[E_k \rightarrow \text{elements of } G \text{ fixing } E_k\]
    and 
    \[\text{fixed field of } H_k \leftarrow H_k.\]
    Under this correspondence:
    \begin{enumerate}
        \item This correspondence is inclusion-reversing, i.e if \(E_1 \to H_1, E_2 \to H_2\), then \(E_1\subseteq E_2 \Leftrightarrow H_2 \leq H_1.\)
        \item \([K:E_k] = H_k\) and \([E_k:F] = \mid G: H_k\mid\)
        \item \(K/E_k\), the intermediate fields over \(F\), is always Galois, with \(Gal(K/E_k) = H_k.\)
        \item \(E_k/F\) is Galois if and only if \(H_k \leq G\) is a normal subgroup.
        \item If \(E_1, E_2\) correspond to \(H_1, H_2\), then \(E_1 \bigcap E_2\) corresponds to \(\rangle H_1, H_2 \langle\) and the composite field 
        \(E_1E_2\) corresponds to \(H_1 \bigcap H_2.\)\\
    \end{enumerate}
    This correspondence between subfields of \(K\) and subgroups of \(Gal(K/F)\) is called the \textit{Galois Correspondence}.
    \end{theorem}
    \textit{Note.} The proof of the Fundamental Theorem of Galois Theory is a must-read for everyone; however, due to length concerns of the main section 
    of the handout, it is in the Proofs section. Please check out some of the other proofs that have faded to obscurity down here.
\end{section}

\begin{section}{Proofs}
    \begin{subsection}{Theorem 4.1: Fundamental Theorem of Galois Theory}
        We will prove each of the 5 statements from FTGT. Some of them requires preliminary results, which we will derive along the way:\\
        \noindent\textbf{(1)} If \(E_1 \mapsto H_1, E_2 \mapsto H_2\), then \(E_1\subseteq E_2\Leftrightarrow H_2 \leq H_1\) (i.e. the correspondence is 
        inclusion-reversing.)
        \begin{proof}
            We've already done this: see \textbf{Proposition 3.1}.
        \end{proof}

        \noindent\textbf{(2)} \([K:E_k] = \left| H_k \right|\) and \([E_k:F] = \left| G:H_k \right|\):
        \begin{figure}[H]
            \centering
            \begin{tikzcd}[row sep=small, arrows=no head]
                K \arrow[d, "\left| H_k \right|"]\\
                E_k\arrow[d, "\left| G:H_k \right|"]\\
                F
            \end{tikzcd}
        \end{figure}
        \begin{proof}
            The proof uses the following theorem:
            \begin{theorem}{(Order of subgroups \& Degree of Extensions)}
                Let \(G\) be a subgroup of \(Aut(K)\). Let \(F\) be the fixed field of \(G\). Then 
                \[[K:F] = \left| G \right|.\]
            \end{theorem}
            Since \(E_k\) is the fixed field of \(H_k\) by the Galois Correspondence, \([K: E_k] = \left| H_k \right|\) holds by \textbf{Theorem 5.1}.
            \(F\) is the fixed field of \(G\), so \([K:F] = \left| G \right|\) also by \textbf{Theorem 5.1}. Also 
            \([K:F] = [K:E_k][E_k:F]\) by the Tower Law, so \([E_k:F] = \frac{\left|G\right|}{\left|H_k\right|} = \left|G:H_k\right|\).
        \end{proof}

        \noindent\textbf{(3)} \(K/E_k\) is always Galois, with \(Gal(K/E_k) = H_k\).
        \begin{proof}
            The proof uses the following corollary:
            \begin{corollary}{(General form of \textbf{Proposition 3.2}.)}
                Let \(K/F\) be any finite extension. Then 
                \[Aut(K/F) \leq [K:F]\]
                with equality if and only if \(F\) is the fixed field of \(Aut(K/F)\).
            \end{corollary}
            By definition of the Galois Correspond,ce \(E_k\) is the fixed field of \(H_k\) and \(H_k = Aut(K/E_k)\). Then 
            \(K/E_k\) is Galois since \(K/E_k = \left|H_k\right|\) (consider the conditions for Galois group and extensions),
            and \(H=Aut(K/E_k) = Gal(K/E_k)\).
        \end{proof}

        \noindent\textbf{(4)} \(E_k/F\) is Galois if and only if \(H_k\leq G\) is a normal subgroup.
        \begin{proof}
            Surprisingly (for now), this is by far the most complicated result of the FTGT to prove and this therefore not included here. 
            I might write a writeup someday for the things that are too hard to be included even in the Proofs sections :(\\
            \noindent\textit{If this is your second or more read-through:} Unsurprisingly, even though this result seems like a variant of \textbf{(2)} except 
            in the opposite direction (we descend towards \(F\) instead of ascending towards \(K\)), this is the by far the most complicated result of 
            the FTGT to prove because it forms the final piece of the `chain of solvable groups' idea in solvability.
        \end{proof}

        \noindent\textbf{(5)} If \(E_1, E_2\) correspond to \(H_1, H_2\), then \(E_1 \cap E_2\) corresponds to \(\langle H_1,H_2 \rangle\) and 
        \(E_1 E_2\) corresponds to \(H_1 \cap H_2\).
        \begin{proof}
            An explanation of what the notation means:
            \begin{enumerate}
                \item \(E_1 \cap E_2\): the intersection of \(E_1\) and \(E_2\), intuitively `\(E_1\) and \(E_2\)'. Note that this is not always the elements that are both in 
                \(E_1\) and \(E_2\), since the resultant set must satisfy the field axioms.
                \item \(\langle H_1,H_2 \rangle\): the smallest subgroup of \(G\) containing all elements in \(H_1, H_2\). Similar to the definition of a
                union of two sets, except in this case extra elements may be added so the resulting set satisfies the group axioms. Intuitively `\(H_1\) or \(H_2\)'.
                \item \(E_1 E_2\): the composite field of \(E_1\) and \(E_2\), essentially the smallest subfield of \(K\) containing all elements in \(E_1\)
                and \(E_2\). Intuition same as above: `\(E_1\) or \(E_2\)'.
                \item: \(H_1 \cap H_2\): the intersection of two groups. You probably already know this as `\(H_1\) or \(H_2\)'.
            \end{enumerate}
            Since \(H_1\) fixes \(E_1\) and \(H_2\) fixes \(E_2\), \(E_1 \cap E_2\) is fixed by all the elements of both \(H_1\) and \(H_2\) (`\(H_1\) or \(H_2\)'), so \(E_1 \cap E_2\) 
            corresponds to \(\langle H_1, H_2 \rangle\). Similarly, \(H_1\cap H_2\) fixes elements in both \(E_1\) and \(E_2\)(`\(E_1\) or \(E_2\)'), so \(E_1 E_2\) corresponds to \(H_1\cap H_2\).\\
            \noindent\textit{Note.} If you ever learned a programming language (that is not HTML, I hate that thing) you'll probably have a good understanding of the `and-or' inverse correspondence 
            coming into play here, similar to `not A and B' corresponds to `not A or not B'.
        \end{proof}
    \end{subsection}
\end{section}
\end{document}